\justifying
\NSFCBodyText

\subsubsection{年度研究计划}
项目按资料建议采用\textbf{【暂定 自正式启动起12个月】}的研究周期。正式起止日期及合同要求以\textbf{【待补 F-KS-05：资助起止日期和结题成果要求】}为准。以下月份为相对启动月份，依据三项研究内容的依赖关系安排，不预先写成已批准的日历日期。

\textbf{M1—M2：建立共同评测与状态标注。}完成官方课题目标和任务边界核对，复现未压缩历史、滑动窗口、固定摘要及轻量压缩基线，建立原文档案和完整成本日志。通过开发任务检查工具接口、环境复位和终态测试，整理动作前提与证据关联标注规则。阶段验收以所有基线能在同一任务清单运行、错误可回溯、开销能核对为标准，并据开发分歧率确定主测试规模。同步按实测GPU时、接口调用数和可分配资源冻结扩展规模；若功效或资源不足，优先完成单环境单模型机制实验，延后额外模型和窗口，不放宽原非劣界。

\textbf{M3—M5：研究任务状态保真压缩。}完成规则与轻量信息价值估计，比较不同表示粒度及依赖保留方式，开展删除、随机保留和关系打乱消融。同步建立相同长度的机制对照和原文回读接口，形成第一版紧凑状态。若轻量学习不能稳定超过规则方案，保留规则作为后续恢复实验的可解释输入，不在这一阶段追求大规模训练。

\textbf{M6—M8：完成有效恢复与核心交叉实验。}实现对象关联、版本核对和必要环境重查，构造有无更新、有无强依赖的成对任务，运行依赖保留与有效性校验的交叉比较。重点核验过期证据导致的真实动作失败及其变化，形成失效案例和人工复核记录。对无法匹配原生模型的近期方法，完成同骨干机制对照并单列原生实现差异。

\textbf{M9—M10：集成调度与外部验证。}接入全开销计量，比较规则与轻量调度，在预先固定的任务划分上报告非劣结果和成本权衡。完成 WebArena 跨环境测试及元数据错误压力测试。企业场景的开展以\textbf{【待补 F-KS-03：业务数据授权、接口与协作负责人】}落实为前提；未落实时提交公开模拟环境结果和迁移要求，不把构造任务标记为真实业务验收。

\textbf{M11—M12：整理成果与复现验收。}冻结可交付版本，核对论文、代码、数据使用许可和第三方依赖，整理全部正负结果及适用边界。由未参与主要方法实现的项目成员复跑核心比较；对无法复现的结果重新定位并报告原因。依据正式合同准备结题材料，最终验收以研究记录、方法实现和可复核证据的一致性为基础。

\subsubsection{预期研究结果}
预期形成\textbf{【暂定 1套ACORE上下文管理原型、1套公开环境评测与错误分析工具、1份可复现技术报告、1—2篇论文投稿】}。原型提供保真压缩、可追溯原文档案、有效恢复与预算调度接口，工具包含任务清单、固定配置、成本统计及错误分型。论文投稿是计划产出，不保证录用；知识产权申请在实际形成可保护成果并核对合作协议后决定，不预列专利号或承诺授权数量。

研究结论包括机制有效与无效的任务条件、不同资源预算下的权衡关系以及可复核负结果。拟公开不受第三方限制的代码、配置和任务描述；受限模型、原始数据及企业资产只提供符合授权范围的复现实验说明。正式成果数量和交付范围依据课题合同终核，研究方法的比较不因成果数量目标而选择性报告。
